\documentclass[10pt,a4paper]{amsart}
\usepackage[margin=19mm]{geometry}
\usepackage{amsmath,amssymb,mathtools}
\usepackage[scr=rsfs,cal=euler]{mathalfa}
\usepackage{xfrac}
\usepackage[unicode,hidelinks,bookmarks=false]{hyperref}
\newcommand{\ie}{i.e.\ }
\newcommand{\cf}{cf.\ }
\newcommand{\resp}{respectively\ }
\newcommand{\Subsection}{\subsection}
\providecommand{\nobreakdash}{\nobreak\hbox{-}}
\setlength{\emergencystretch}{2em}
\multlinegap0pt
\usepackage{bm}
\usepackage{bbm}
\usepackage{dsfont}
\usepackage{xspace}
\usepackage{cancel}
\usepackage{mathtools}
\usepackage{amsthm}
\makeatletter
\@ifundefined{thm}{\newtheorem{thm}{Thm}}{}
\@ifundefined{lemma}{\newtheorem{lemma}[thm]{Lemma}}{}
\@ifundefined{prop}{\newtheorem{prop}[thm]{Proposition}}{}
\makeatother
\makeatletter
\newcommand{\Leb}{\operatorname{Leb}}
\newcommand{\cD}{{\mathcal D}}
\newcommand{\eps}{\varepsilon}
\newcommand{\hcD}{{\widehat {\mathcal D}}}
\expandafter\ifx\csname pfof\endcsname\relax
\newenvironment{pfof}[1]{\vspace{1ex}\noindent{\bf Proof of
#1}\hspace{0.5em}}{\hfill\qed\vspace{1ex}}
\fi
\newcommand{\tcD}{{\widetilde {\mathcal D}}}
\makeatother
\title[Good inducing schemes for uniformly hyperbolic flows, and ap]{Good inducing schemes for uniformly hyperbolic flows, and applications to exponential decay of correlations}
\author{Ian Melbourne, Paulo Varandas}
\date{Source: arXiv:2302.12363v3 (CC BY 4.0)}
\begin{document}
\maketitle

\noindent\emph{Source and licence.} Good inducing schemes for uniformly hyperbolic flows, and applications to exponential decay of correlations. Version arXiv:2302.12363v3 (CC BY 4.0), distributed under the Creative Commons Attribution 4.0 International licence, as recorded in the arXiv licence metadata for this exact version and shown on the abstract page https://arxiv.org/abs/2302.12363v3. Full rights notice: licenses/arxiv-2302.12363-CC-BY-4.0.txt.\par\smallskip

\noindent\textit{Editorial scope.} Counted unit: the Lemma whose source label is \texttt{lem:keyfact} in the article (source0.tex lines 607--614, source label \texttt{lem:keyfact}), delivered together with the article's own complete original proof of it (source0.tex lines 616--760, one \texttt{proof} environment of 145 lines). No mathematics is written, repaired or re-derived: statement and proof are the article's own text line for line. Required context retained uncounted: eq:bdd (lines 240--242); eq:bddpi (lines 247--249); prop:eps (lines 330--333); prop:eps2 (lines 400--409). Every definition, display and label of those lines is retained unchanged. Omitted from this package and not redistributed: 1--239, 243--246, 250--329, 334--399, 410--606, 615--615, 761--1663. Nothing inside a retained excerpt is omitted or altered: every retained body is delivered exactly as the article's lines are written, so no omission receipt of any kind accompanies this package. Citations: the retained excerpts carry no citation command at all, so no bibliography entry of the article is transcribed and no reference is redistributed. Reference adaptation: none was needed and none was made; every reference inside the delivered unit resolves inside it, so no unresolved reference or citation is produced. Printed slips kept literal: no printed word, symbol, hypothesis or number of the counted statement or of its proof was altered. Layout and class: the article's own class is \texttt{article} with packages including graphicx, amsmath, amssymb, theorem, hyperref; the wrapper is the lane's pinned \texttt{amsart} editorial wrapper at 10pt on A4, which restates the theorem-like environments the retained text needs and adds the article's own macro definitions that the retained text needs (\texttt{\textbackslash Leb}, \texttt{\textbackslash cD}, \texttt{\textbackslash eps}, \texttt{\textbackslash hcD}, \texttt{\textbackslash tcD}), with extra wrapper packages (bm, bbm, dsfont, xspace, cancel, mathtools). The delivered numbering is the unit's own. Assets: no retained span contains a figure environment, a graphics-inclusion command, a PDF-page inclusion, a table or a TikZ picture; no image file is copied into this package, the package declares no asset dependency and no image file is redistributed with it. Licence: version arXiv:2302.12363v3 of this article is distributed under the Creative Commons Attribution 4.0 International licence, as recorded in the arXiv licence metadata for this exact version and shown on the abstract page https://arxiv.org/abs/2302.12363v3. A full rights notice is delivered at licenses/arxiv-2302.12363-CC-BY-4.0.txt. Characters: every retained excerpt is pure ASCII, so the delivered engine, XeLaTeX with its default Latin Modern text font, reports no missing character.
\begin{equation}\label{eq:bdd}
\frac{|\det D\phi_n(x)|E^u|}{|\det D\phi_n(y)|E^u|} \leq C_1
\end{equation}

\begin{equation}\label{eq:bddpi}
C_2^{-1}\le \frac{\Leb(\pi(E))}{\Leb(E)} \leq C_2
\end{equation}

\begin{prop} \label{prop:eps}
Choose $\eps<(C_3^{-1}\delta)^{1/\alpha}$ sufficiently small that $W^u_\eps(x)\subset\hcD$ for all $x\in\tcD_L$. Then
 $\bigcup_jU_{nj}^{L-1}\subset A_{n-1}$ for all $n\ge1$.
\end{prop}

\begin{prop} \label{prop:eps2}
Choose $\eps<\big\{C_3^{-1}\delta(\lambda^{-\alpha}-1)\big\}^{1/\alpha}$.
Then for all $n\ge1$,
\begin{itemize}
\item[(a)]
$A_{n-1}^{(\eps)}\subset \{y\in Y_{n-1}:t_{n-1}(y)\le 1\}$ for all $n\ge1$.
\item[(b)]
$\phi_{-n}W_\eps^u(\phi_nx)\subset A_{n-1}^{(\eps)}$ for all $x\in A_{n-1}$.
\end{itemize}
\end{prop}

\begin{lemma} \label{lem:keyfact}
Let $\eps\in(0,\frac12\delta_0)$ be small as in Propositions~\ref{prop:eps} and~\ref{prop:eps2}.
There exist $c_1>0$ and $N\ge1$ such that 
\[
\Leb\Big(\bigcup_{i=0}^{N} \{R=n+i\}\Big)\ge c_1 \Leb(A_{n-1})
\quad\text{for all $n\ge1$}.
\]
\end{lemma}

\begin{proof}
Fix $\lambda\in(0,1)$, $L>1$, $0<\delta<\delta_1<\delta_0$ and $N_1\ge1$ as defined from the outset.
Recall that $C_3(3\delta)^\alpha<\frac12\delta_0$ and $C_4(L+1)\delta<\delta_0$.
Choose $N_2\ge1$ such that $\lambda^{N_2}<\eps/\delta_0$ and take $N=N_1+N_2$.

We claim that
\begin{itemize}
\item[(*)]
For all $z\in \Lambda$, there exists $i\in\{1,\dots,N_1\}$ such that
$\pi (\phi_{i+N_2} W^u_\eps(z)\cap \tcD_L)\supset \cD_L$.
\end{itemize}

Fix $z\in \Lambda$.
By the definition of $N_1$, there exists $1\le i\le N_1$ such that $d_M(\phi_{-i}p,\phi_{N_2}z)<\delta$.
Let $y\in \cD_L$. Then
\[
d_M(\phi_{-i}y,\phi_{N_2}z)\le 
d(\phi_{-i}y,\phi_{-i}p)+d_M(\phi_{-i}p,\phi_{N_2}z)\le d(y,p)+d_M(\phi_{-i}p,\phi_{N_2}z)<(L+1)\delta.
\]
Using the local product structure and choice of $\delta$, we can
define $x\in W^{cs}_{\delta_1}(\phi_{-i}y)\cap W^u_{\delta_1}(\phi_{N_2}z)$.
Then $\phi_ix\in W^{cs}_{\delta_1}(y)\subset\tcD_L$ and
$g_ix=\pi \phi_ix=y$.
Also,
\[
d(x, \phi_{N_2}z)\le C_4d_M(\phi_{-i}y,\phi_{N_2}z)
<C_4(L+1)\delta<\delta_0.
\]
By the definition of $N_2$,
\[
\phi_ix\in \phi_i W^u_{\delta_0}(\phi_{N_2}z)
\subset \phi_{i+N_2} W^u_\eps(z).
\]
Hence we obtain that
$y=\pi \phi_ix\in \pi(\phi_{i+N_2}W^u_\eps(z)\cap\hcD_L)$
proving (*).

Next, we claim that 
\begin{itemize}
\item[(**)] 
For all $z\in \phi_nA_{n-1}$, $n\ge1$, there 
exist $i\in\{0,\dots,N\}$ and $j$ such that $U_{n+i,j}^1\subset \phi_{-n}W^u_{\delta_0}(z)$.
\end{itemize}

To prove (**), define $V_\eps=\phi_{-n}W^u_\eps(z)$.
By Proposition~\ref{prop:eps2}(b),
$V_\eps\subset A_{n-1}^{(\eps)}$.
We now consider two possible cases.

Suppose first that $V_\eps\subset A_{n+i}$ for all $0\le i\le N$.
By claim (*), there exists $1\le i\le N=N_1+N_2$ such that 
\[
\pi(\phi_{n+i}V_\eps\cap \tcD_L)=\pi(\phi_iW^u_\eps(z)\cap \tcD_L)\supset \cD_L,
\]
while $V_\eps\subset A_{n+i-1}$ by assumption.
This means that $V_\eps\supset U_{n+i,j}^L$ for some $j$. Hence
\[
U_{n+i,j}^1\subset U_{n+i,j}^L\subset V_\eps\subset  \phi_{-n}W^u_{\delta_0}(z),
\]
  and we are done.

In this way, we reduce to the second case where there exists $0\le i\le N$ least such that
$V_\eps\not\subset A_{n+i}$.
Since $i$ is least, 
$V_\eps\subset A_{n+i-1}^{(\eps)}$. (The $\eps$ is required in case $i=0$.)
By Proposition~\ref{prop:eps2}(a),
$V_\eps\subset\{t_{n+i-1}\le 1\}$.
Hence
\[
\begin{split}
V_\eps\setminus A_{n+i}
& =(V_\eps\cap B_{n+i})\,\cup\,(V_\eps\cap \{R=n+i\})
\\ & \subset\{t_{n+i-1}\le 1,\,t_{n+i}\ge1\}\,\cup\, \{R=n+i\}
\subset\bigcup_j U_{n+i,j}^2.
\end{split}
\]
Since $V_\eps\setminus A_{n+i}\neq\emptyset$, 
this means that there exists $j$ so that
$V_\eps$ intersects $U_{n+i,j}^2$.
Hence we can choose $a_2\in W^u_\eps(z)\cap \phi_n U_{n+i,j}^2$.

Recall that $\phi_{n+i}U_{n+i,j}^m \subset \tcD_m$ 
and $g_{n+i}U_{n+i,j}^m = \cD_m$ for $m=1,2$.
In particular, $b_2=\phi_ia_2\in \tcD_2$ and $c_2=g_ia_2\in \cD_2$.

Let $c_1\in\cD_1$.
Then $d_M(c_1,b_2)\le d_M(c_1,c_2)+d_M(c_2,b_2)<3\delta+\delta_1<4\delta_1$. Hence,
using the local product structure and definition of $\delta_1$,  we can define $b_1\in W^{cs}_{\delta_0}(c_1)\cap W^u_{\delta_0}(b_2)$ and $a_1=\phi_{-i}b_1$.
Note that
\[
\phi_ia_r=b_r,\qquad \pi b_r=c_r, \quad r=1,2.
\]
Hence
\[
d(a_1,a_2)\le d(b_1,b_2)\le C_3d(c_1,c_2)^\alpha<C_3(3\delta)^\alpha<\tfrac12\delta_0,
\]
and so $d(a_1,z)\le d(a_1,a_2)+d(a_2,z)<\frac12\delta_0+\eps<\delta_0$.
It follows that
$a_1\in W^u_{\delta_0}(z)$ and thereby that
$c_1\in g_i(W^u_{\delta_0}(z)\cap \phi_{-i}\hcD_1)$.
This proves that $\cD_1\subset g_i(W^u_{\delta_0}(z)\cap \phi_{-i}\hcD_1)$.
Hence $U_{n+i,j}^1\subset g_{-(n+i)}\cD_1\subset \phi_{-n}W^u_{\delta_0}(z)$ verifying claim~(**).

 \vspace{1ex}
We are now in a position to complete the proof of the lemma.
Let $n\ge1$, and let $Z\subset \phi_nA_{n-1}$ be a maximal set of points such that
the balls $W^u_{\delta_0/2}(z)$ are disjoint for $z\in Z$.
If $x\in \phi_nA_{n-1}$, then $W^u_{\delta_0/2}(x)$ intersects at least one
$W^u_{\delta_0/2}(z)$, $z\in Z$, by maximality of the set $Z$. Hence
$\phi_nA_{n-1}\subset \bigcup_{z\in Z}W^u_{\delta_0}(z)$.
It follows that
\[
A_{n-1}\subset \bigcup_{z\in Z}\phi_{-n}W^u_{\delta_0}(z).
\]

Let $z\in Z$ and let $U_z=U_{n+i,j}^1$ be as in claim~(**).
In particular, $g_{n+i}U_z=\cD_1=W^u_\delta(p)$.
Also, $\Leb(\phi_{n+i}U_z)\le 
|D\phi_1|_\infty^{im}\Leb(\phi_nU_z)$
where $m=\dim E^u$.
Hence, by~\eqref{eq:bddpi}, 
\[
\frac{1}{\Leb(\phi_nU_z)}\le |D\phi_1|_\infty^{Nm}\frac{1}{\Leb(\phi_{n+i}U_z)}
\le C_3 |D\phi_1|_\infty^{Nm}\frac{1}{\Leb(W^u_\delta(p))}.
\]
By~\eqref{eq:bdd}, 
\[
\frac{\Leb(\phi_{-n}W^u_{\delta_0}(z))}{\Leb(U_z)}
\le C_1
\frac{\Leb(W^u_{\delta_0}(z))}{\Leb(\phi_nU_z)}
\le 
K,
\]
where $K=C_1C_3|D\phi_1|_\infty^{Nm}\,\dfrac{\sup_{y\in Y}\Leb(W^u_{\delta_0}(y))}{\Leb(W^u_\delta(p))}$.

Finally, the sets $U_z$ are connected components of $\bigcup_{0\le i\le N}\{R=n+i\}$ lying in distinct disjoint sets $\phi_{-n}W^u_{\delta_0}(z)$.
Hence 
\[
\Leb(A_{n-1})  
 \le \sum_{z\in Z} \Leb(\phi_{-n}W^u_{\delta_0}(z))
 \le K \sum_{z\in Z}\Leb(U_z) \le 
K\Leb\Big(\bigcup_{0\le i\le N}\{R=n+i\}\Big),
\]
as required.
\end{proof}

\end{document}
